
\subsubsection{5.1 Quality Control: Mock Data Detection and Correction}

Automated validator detected five synthetic data patterns in proof-of-concept code (main_poc.py lines 50, 60, 65, 67, and logic analysis 60-67). Violations included: np.random.uniform(0.75, 0.90) generating synthetic baseline accuracies, base_accuracy = subject_difficulties[subject] ensuring same accuracy across all λ by design, np.random.normal(0, 0.05) adding random noise instead of real data computation, np.random.random() < prob_correct generating boolean results from random process, and tautological logic guaranteeing ICC > 0.95 by construction.

Corrective actions (Attempt 1/5): disabled main_poc.py (renamed to main_poc.py.DISABLED), verified main.py uses real data loaders (MMLULoader, HumanEvalLoader), confirmed MMLULoader loads from HuggingFace cais/mmlu dataset, confirmed HumanEvalLoader loads from OpenAI human-eval pip package, cleaned old POC synthetic results, verified real data loading works correctly.

Fix verification passed 6/6 checks: POC code disabled, main.py active, MMLU 57 subjects loaded, HumanEval 164 problems loaded, no synthetic patterns in main.py, datasets match published benchmarks.

This demonstrates automated quality control effectiveness. The validator identified that high ICC > 0.95 was guaranteed by synthetic data construction, not genuine capability invariance. Without this detection, h-e1 would have falsely validated the prerequisite assumption, allowing invalid experiments to proceed.

\subsubsection{5.2 Dataset Verification}

\textbf{MMLU Dataset:}
\begin{itemize}
\item Source: HuggingFace cais/mmlu
\item Subjects: 57 (abstract_algebra through virology)
\item Test Questions: 14,042
\item Question Format: Multiple choice, 4 options
\item Few-Shot: 4 examples per subject
\item Sample Verified: "This question refers to the following information. 'The far-reaching, the boundless future will be the era of American greatness...' By what means did the United States take possession of the Oregon Territory?" (high_school_us_history)
\end{itemize}

\textbf{HumanEval Dataset:}
\begin{itemize}
\item Source: OpenAI human-eval (pip package)
\item Problems: 164 hand-written programming challenges
\item Format: Function signature + docstring
\item Evaluation: Execution-based correctness (pass@1)
\item Sample Verified: HumanEval/0: has_close_elements - "Check if in given list of numbers, are any two numbers closer to each other than given threshold"
\end{itemize}

Dataset verification confirms infrastructure correctness. The experiment would evaluate on standard, widely-used benchmarks, ensuring comparability with prior work and reproducibility.

\subsubsection{5.3 Verification Protocol Design}

\textbf{Hypothesis Chain:}

\begin{table}[htbp]
\centering
\begin{tabular}{llllll}
\toprule
Hypothesis & Type & Gate & Prerequisite & Success Criterion & If Fail \\
\midrule
h-e1 & EXISTENCE & MUST_WORK & None & (ICC > 0.95) AND (p > 0.05) AND (f < 0.10) & ABORT chain \\
h-m1 & MECHANISM & DETERMINES_SUCCESS & h-e1 PASS & (R² ≥ 0.6) AND (ΔAIC > 10) & MODIFY hypothesis \\
h-m2 & MECHANISM & MUST_WORK & h-m1 PASS & r > 0.5, p < 0.05 & ABORT chain \\
h-m3 & MECHANISM & DETERMINES_SUCCESS & h-m2 PASS & HOR difference ≥ 10\%, p < 0.05 & MODIFY hypothesis \\
h-m4 & MECHANISM & DETERMINES_SUCCESS & h-m3 PASS & k₂_ratio ≥ 1.30, 95\% CI excludes 1.0 & MODIFY hypothesis \\
\bottomrule
\end{tabular}
\end{table}

\textbf{Falsification Criteria:}

\begin{table}[htbp]
\centering
\begin{tabular}{lll}
\toprule
Assumption & Falsifier & Consequence if Violated \\
\midrule
A1: Capability invariance & ICC < 0.95 OR p < 0.05 & ACE confounded with capability \\
A2: Error validity & Real errors ≠ confusion matrix & HOR measures stress-testing \\
A3: Automation bias & Mediation test fails & Alternative mechanism \\
A4: Temporal dynamics & Neither model fits (R² < 0.6) & Alternative patterns exist \\
A5: Intervention practicality & Frustration ≥ 7/10 OR abandonment ≥ 20\% & Intervention impractical \\
\bottomrule
\end{tabular}
\end{table}

This demonstrates pre-specification of success criteria, failure handling, and iterative refinement protocols. Dependency structure ensures validity checks are enforced (h-e1 must pass before h-m1 can test coupling).

\subsubsection{5.4 Deviation Analysis}

\textbf{Planned vs. Actual:}

\begin{table}[htbp]
\centering
\begin{tabular}{lllll}
\toprule
Hypothesis & Planned Metric & Planned Target & Actual Result & Deviation Type \\
\midrule
h-e1 & ICC & ICC > 0.95 & NOT COMPUTED & IMPLEMENTATION_GAP \\
h-e1 & ANOVA p-value & p > 0.05 & NOT COMPUTED & IMPLEMENTATION_GAP \\
h-e1 & Cohen's f & f < 0.10 & NOT COMPUTED & IMPLEMENTATION_GAP \\
\bottomrule
\end{tabular}
\end{table}

All deviations are IMPLEMENTATION_GAP (infrastructure/resource constraint), not HYPOTHESIS_ISSUE (scientific flaw) or DESIGN_ISSUE (experimental design problem). The experiment design and implementation were completed successfully. The gap is execution resources (API key, budget approval approximately $1,620 for full MMLU + HumanEval evaluation across 5 λ levels with approximately 4 hours runtime requirement).
